\section{Esercizio B -- Sistema LTI a tempo discreto}

Il sistema è $x(k+1)=Ax(k)+Bu(k)$, $y(k)=Cx(k)$, SISO e proprio, con
\[
A=\begin{bmatrix}
-\tfrac{92}{125} & -\tfrac{107}{125} & -\tfrac{43}{50}\\[2pt]
\tfrac{217}{125} & -\tfrac{18}{125} & -\tfrac{7}{50}\\[2pt]
-\tfrac{184}{125} & \tfrac{36}{125} & \tfrac{7}{25}
\end{bmatrix},\quad
B=\begin{bmatrix}-\tfrac12\\[2pt]\tfrac12\\[2pt]-1\end{bmatrix},\quad
C=\begin{bmatrix}\tfrac94 & -\tfrac74 & -2\end{bmatrix}.
\]

\subsection{Modi naturali}
Nei sistemi a tempo discreto i modi naturali sono successioni reali nella forma
$\lambda_i^{\,k}$. La matrice $A$ ha l'unico autovalore $\lambda=-\tfrac15$ con
molteplicità algebrica $3$ ma geometrica $1$: $A$ \emph{non} è diagonalizzabile.
\texttt{JordanDecomposition} restituisce la matrice di cambio base $T$ e la forma
canonica $\Lambda$ con un solo blocco di Jordan $3\times3$. I tre autovalori
coincidenti generano modi di tipo binomiale, secondo
\[
\sum_{i=0}^{l-1}\binom{k}{i}\lambda^{\,k-i},
\]
cioè $\lambda^k$, $\binom{k}{1}\lambda^{k-1}$ e $\binom{k}{2}\lambda^{k-2}$.

\mmacap[0.48]{img-084}{1\textordmasculine\ modo: $(-\tfrac15)^k$}
\mmacap[0.48]{img-085}{2\textordmasculine\ modo: $\binom{k}{1}(-\tfrac15)^{k-1}$}
\mmacap[0.48]{img-086}{3\textordmasculine\ modo: $\binom{k}{2}(-\tfrac15)^{k-2}$}

\subsection{Risposta libera}
La risposta libera sullo stato è $x(k)=A^k x_0$. Analogamente al caso continuo si cerca
una matrice simile ad $A$ tramite $T$: con $\mathrm{MatrixPower}[\Lambda,k]$ si ottiene
una matrice con, sulla diagonale, le potenze $k$-esime dei modi e, sulla triangolare
superiore, le successioni polinomio-potenza. Quindi $x_l(k)=T\,\Lambda^k z_0$ con
$z_0=T^{-1}x_0$, e $y_l(k)=C\,x_l(k)$. Lo stato iniziale è
$x_0=\begin{bmatrix}0&0&-3\end{bmatrix}^{\!\top}$.

\mmacap[0.48]{img-091}{Risposta libera sull'uscita $y_l[k]$}
\mmacap[0.48]{img-092}{Risposta libera sullo stato e sull'uscita a confronto}

\subsection{Attivazione di specifici modi naturali}
Come nel caso continuo, scegliendo opportune combinazioni lineari delle colonne di $T$
si attivano sulla risposta libera soltanto i modi desiderati.

\mmacap[0.48]{img-094}{Modi legati alle ultime due colonne di $T$}
\mmacap[0.48]{img-096}{Modi legati alla prima e all'ultima colonna}
\mmacap[0.48]{img-098}{Solo prima colonna}
\mmacap[0.48]{img-100}{Solo seconda colonna}
\mmacap[0.5]{img-102}{Solo terza colonna}

\subsection{Funzione di trasferimento, poli e zeri}
La funzione di trasferimento è ora funzione della variabile complessa $z$:
$Y(z)=G(z)U(z)$ con $G(z)=C(zI-A)^{-1}B+D$. Calcolata dalla definizione e normalizzata,
coincide con quella prodotta da \texttt{StateSpaceModel}. I poli sono le radici del
denominatore e, come sempre, sono un sottoinsieme degli autovalori: qui il polo triplo
$z=-\tfrac15$. Lo zero (radice del numeratore) è $z=-\tfrac18$.

\subsection{Risposta al gradino unitario}
La risposta forzata al gradino è $Y_f(z)=G(z)\,U(z)$ con $\Ztr[1(k)]=\tfrac{z}{z-1}$.
Si divide $Y_f(z)$ per $z$ per operare come nel tempo continuo e si scrivono i fratti
semplici. Per il residuo legato al polo $z=1$ (componente di regime) basta la formula di
Heaviside semplificata; per il polo triplo $z=-\tfrac15$ occorre quella generalizzata,
con le derivate successive
\[
C_{2j}=\frac{1}{(3-j)!}\lim_{z\to-\frac15}\frac{d^{\,3-j}}{dz^{\,3-j}}
\Big[(z+\tfrac15)^3\tfrac{Y_f(z)}{z}\Big].
\]
Antitrasformando (i modi binomiali danno i termini $\binom{k}{j}(-\tfrac15)^{k-j}$) si
ottiene la risposta come somma di componente di regime e transitoria.

\mmacap[0.7]{img-115}{Risposta al gradino con componenti $y_f[k]$, $y_{ss}[k]$, $y_{tr}[k]$}

Si verifica che la differenza poli$-$zeri (pari a $2$) è rispettata.

\subsection{Modello ARMA}
Partendo da $G(z)=\tfrac{Y(z)}{U(z)}$ e dalla matrice di osservabilità si ricavano le
condizioni iniziali sull'uscita compatibili con
$x_0=\begin{bmatrix}0&-2&-3\end{bmatrix}^{\!\top}$. Uguagliando
$\mathrm{den}(G)\,Y(z)=\mathrm{num}(G)\,U(z)$ e antitrasformando $z^nY(z)\to y(k+n)$ si
ottiene l'equazione alle differenze; nel dominio $z$ si calcola la risposta del sistema
e se ne isola la componente di risposta libera.

\mma[0.85]{img-122}

Confrontando la risposta libera dal modello ARMA con quella dalla definizione si verifica
l'uguaglianza, a conferma della correttezza delle condizioni iniziali e quindi
dell'equivalenza delle due rappresentazioni.

\paragraph{Risposta all'ingresso a finestra.}
L'ingresso vale $1$ per $0\le k\le10$ e $0$ altrove. Essendo di durata finita, la sua
trasformata $\Ztr$ si scrive come somma di potenze $\sum_{k=0}^{10}z^{-k}$.
Sostituendo questa trasformata e le condizioni iniziali nel modello I/U e
antitrasformando si ottiene la risposta del sistema.

\mmacap[0.48]{img-130}{Risposta totale all'ingresso a finestra}
\mmacap[0.48]{img-132}{Componente di regime della risposta}

\subsection{Condizioni iniziali ARMA che annullano il transitorio}
Con lo stesso ragionamento del caso continuo, si cercano le condizioni iniziali sull'uscita
tali che risposta libera $+$ transitoria sia identicamente nulla. Imponendo l'annullamento
dei coefficienti del numeratore (comandi \texttt{Solve} e \texttt{CoefficientList}) si
ottiene $y_0$, e da questo lo stato iniziale $x_0=\mathrm{Ob}^{-1}y_0$ con
$y_0=\begin{bmatrix}\tfrac{125}{96}&\tfrac{125}{96}&-\tfrac{67}{96}\end{bmatrix}^{\!\top}$.
